% GENERATED FILE. Edit canonical lessons, then run npm run generate:books.
% canonical-source-hash: fnv1a64:2ed5670d312be017
% canonical-lessons: HI-C530-sangit, HI-C530-sur, HI-C530-tal, HI-C530-rag, HI-C530-gungunana

\chapter{Learning a Raga}
\label{ch:hi-learning-a-raga}
\hlchaptermodality{530}

\begin{chapteropening}
\textbf{By the end of this chapter:} \emph{I can talk about music lessons.}

Complete the last lesson of chapter 530: I can talk about music lessons.
\end{chapteropening}

\section[\texorpdfstring{\hi{संगीत} (saṅgīt)}{संगीत (saṅgīt)}]{\hi{संगीत} (saṅgīt) — music}

\label{lesson:HI-C530-sangit}



\begin{quote}
Before the new one: say the Hindi for a character, then the Hindi for a spectator.
\end{quote}

\subsection*{You'll want to know: \hi{संगीत}}
\textbf{\hi{संगीत}} — \emph{saṅgīt} — \textquotedblleft{}music\textquotedblright{}.

\begin{quote}
\hi{मेरी} \hi{बहन} \hi{संगीत} \hi{सीख} \hi{रही} \hi{है।}
\end{quote}

\begin{quote}
\emph{Merī bahan saṅgīt sīkh rahī hai.}
\end{quote}

\begin{quote}
My sister is learning music.
\end{quote}

Masculine.

\begin{cousinweb}
\textbf{\hi{संगीत}} (\emph{saṅgīt}) is a tatsama, taken whole from Sanskrit \emph{saṃgīta}, \textquotedblleft{}sung together; music\textquotedblright{} (Monier-Williams).
\end{cousinweb}

\begin{grammarlens}[title={What you've built}]
One more word for describing what you lived and felt.
\end{grammarlens}

\subsection*{Your turn}
\begin{itemize}
\raggedright
  \item \emph{Say it:} \emph{saṅgīt}
  \item \emph{Say it:} \emph{saṅgīt}, once more
  \item \emph{Say it:} \emph{Merī bahan saṅgīt sīkh rahī hai.} — My sister is learning music.
  \item \emph{Recall:} say the Hindi for a character, then the Hindi for a spectator, then say \emph{saṅgīt} again
\end{itemize}

\begin{tcolorbox}[breakable,colback=teal!4,colframe=teal!35!black,title={Before you move on}]
What does \hi{संगीत} mean? (\textquotedblleft{}Music\textquotedblright{}.) Say it once more.
\end{tcolorbox}
\section[\texorpdfstring{\hi{सुर} (sur)}{सुर (sur)}]{\hi{सुर} (sur) — a musical note; a pitch}

\label{lesson:HI-C530-sur}



\begin{quote}
Before the new one: say the Hindi for a spectator, then the Hindi for music.
\end{quote}

\subsection*{You'll want to know: \hi{सुर}}
\textbf{\hi{सुर}} — \emph{sur} — \textquotedblleft{}a musical note; a pitch\textquotedblright{}.

\begin{quote}
\hi{गायक} \hi{का} \hi{सुर} \hi{बहुत} \hi{मीठा} \hi{है।}
\end{quote}

\begin{quote}
\emph{Gāyak kā sur bahut mīṭhā hai.}
\end{quote}

\begin{quote}
The singer's voice is very sweet.
\end{quote}

Masculine.

\begin{cousinweb}
\textbf{\hi{सुर}} (\emph{sur}) is a tadbhava: CDIAL lists Hindi \emph{sur}, \textquotedblleft{}sound, words, tune\textquotedblright{}, under Sanskrit \emph{sura}, \textquotedblleft{}sound\textquotedblright{}, a variant of \emph{svara}, \textquotedblleft{}a musical note\textquotedblright{} (CDIAL 13498).
\end{cousinweb}

\begin{grammarlens}[title={What you've built}]
One more word for describing what you lived and felt.
\end{grammarlens}

\subsection*{Your turn}
\begin{itemize}
\raggedright
  \item \emph{Say it:} \emph{sur}
  \item \emph{Say it:} \emph{sur}, once more
  \item \emph{Say it:} \emph{Gāyak kā sur bahut mīṭhā hai.} — The singer's voice is very sweet.
  \item \emph{Recall:} say the Hindi for a spectator, then the Hindi for music, then say \emph{sur} again
\end{itemize}

\begin{tcolorbox}[breakable,colback=teal!4,colframe=teal!35!black,title={Before you move on}]
What does \hi{सुर} mean? (\textquotedblleft{}A musical note; a pitch\textquotedblright{}.) Say it once more.
\end{tcolorbox}
\section[\texorpdfstring{\hi{ताल} (tāl)}{ताल (tāl)}]{\hi{ताल} (tāl) — a beat; a rhythmic cycle}

\label{lesson:HI-C530-tal}



\begin{quote}
Before the new one: say the Hindi for music, then the Hindi for a musical note.
\end{quote}

\subsection*{You'll want to know: \hi{ताल}}
\textbf{\hi{ताल}} — \emph{tāl} — \textquotedblleft{}a beat; a rhythmic cycle\textquotedblright{}.

\begin{quote}
\hi{मीरा} \hi{ताल} \hi{के} \hi{साथ} \hi{हाथ} \hi{हिला} \hi{रही} \hi{है।}
\end{quote}

\begin{quote}
\emph{Mīrā tāl ke sāth hāth hilā rahī hai.}
\end{quote}

\begin{quote}
Mira is moving her hand with the beat.
\end{quote}

Masculine.

\begin{cousinweb}
\textbf{\hi{ताल}} (\emph{tāl}) is a tatsama, taken whole from Sanskrit \emph{tāla}, \textquotedblleft{}clapping the hands; musical time or measure\textquotedblright{} (Monier-Williams).
\end{cousinweb}

\begin{grammarlens}[title={What you've built}]
One more word for describing what you lived and felt.
\end{grammarlens}

\subsection*{Your turn}
\begin{itemize}
\raggedright
  \item \emph{Say it:} \emph{tāl}
  \item \emph{Say it:} \emph{tāl}, once more
  \item \emph{Say it:} \emph{Mīrā tāl ke sāth hāth hilā rahī hai.} — Mira is moving her hand with the beat.
  \item \emph{Recall:} say the Hindi for music, then the Hindi for a musical note, then say \emph{tāl} again
\end{itemize}

\begin{tcolorbox}[breakable,colback=teal!4,colframe=teal!35!black,title={Before you move on}]
What does \hi{ताल} mean? (\textquotedblleft{}A beat; a rhythmic cycle\textquotedblright{}.) Say it once more.
\end{tcolorbox}
\section[\texorpdfstring{\hi{राग} (rāg)}{राग (rāg)}]{\hi{राग} (rāg) — a raga; a melodic mode}

\label{lesson:HI-C530-rag}



\begin{quote}
Before the new one: say the Hindi for a musical note, then the Hindi for a beat.
\end{quote}

\subsection*{You'll want to know: \hi{राग}}
\textbf{\hi{राग}} — \emph{rāg} — \textquotedblleft{}a raga; a melodic mode\textquotedblright{}.

\begin{quote}
\hi{सुबह} \hi{के} \hi{लिए} \hi{अलग} \hi{राग} \hi{होते} \hi{हैं।}
\end{quote}

\begin{quote}
\emph{Subah ke liye alag rāg hote haiṁ.}
\end{quote}

\begin{quote}
There are separate ragas for the morning.
\end{quote}

Masculine.

\begin{cousinweb}
\textbf{\hi{राग}} (\emph{rāg}) is a tatsama, taken whole from Sanskrit \emph{rāga}, \textquotedblleft{}colouring; a musical mode\textquotedblright{} (Monier-Williams).
\end{cousinweb}

\begin{grammarlens}[title={What you've built}]
One more word for describing what you lived and felt.
\end{grammarlens}

\subsection*{Your turn}
\begin{itemize}
\raggedright
  \item \emph{Say it:} \emph{rāg}
  \item \emph{Say it:} \emph{rāg}, once more
  \item \emph{Say it:} \emph{Subah ke liye alag rāg hote haiṁ.} — There are separate ragas for the morning.
  \item \emph{Recall:} say the Hindi for a musical note, then the Hindi for a beat, then say \emph{rāg} again
\end{itemize}

\begin{tcolorbox}[breakable,colback=teal!4,colframe=teal!35!black,title={Before you move on}]
What does \hi{राग} mean? (\textquotedblleft{}A raga; a melodic mode\textquotedblright{}.) Say it once more.
\end{tcolorbox}
\section[\texorpdfstring{\hi{गुनगुनाना} (gungunānā)}{गुनगुनाना (gungunānā)}]{\hi{गुनगुनाना} (gungunānā) — to hum}

\label{lesson:HI-C530-gungunana}



\begin{quote}
Before the new one: say the Hindi for a beat, then the Hindi for a raga.
\end{quote}

\subsection*{You'll want to know: \hi{गुनगुनाना}}
\textbf{\hi{गुनगुनाना}} — \emph{gungunānā} — \textquotedblleft{}to hum\textquotedblright{}.

\begin{quote}
\hi{मीरा} \hi{रसोई} \hi{में} \hi{एक} \hi{पुराना} \hi{गाना} \hi{गुनगुना} \hi{रही} \hi{है।}
\end{quote}

\begin{quote}
\emph{Mīrā rasoī meṁ ek purānā gānā gungunā rahī hai.}
\end{quote}

\begin{quote}
Mira is humming an old song in the kitchen.
\end{quote}

Past, with \hi{ने} (he or she …): \hi{उसने} \hi{एक} \hi{गाना} \hi{गुनगुनाया}, \emph{usne ek gānā gungunāyā}.

\begin{cousinweb}
\textbf{\hi{गुनगुनाना}} (\emph{gungunānā}) is a tadbhava, inherited from \emph{*gunaguna}, \textquotedblleft{}murmur\textquotedblright{}, a form CDIAL reconstructs from its descendants (CDIAL 4013).
\end{cousinweb}

\begin{grammarlens}[title={What you've built}]
One more word for describing what you lived and felt.
\end{grammarlens}

\subsection*{Your turn}
\begin{itemize}
\raggedright
  \item \emph{Say it:} \emph{gungunānā}
  \item \emph{Say it:} \emph{gungunānā}, once more
  \item \emph{Say it:} \emph{Mīrā rasoī meṁ ek purānā gānā gungunā rahī hai.} — Mira is humming an old song in the kitchen.
  \item \emph{Recall:} say the Hindi for a beat, then the Hindi for a raga, then say \emph{gungunānā} again
\end{itemize}

\begin{tcolorbox}[breakable,colback=teal!4,colframe=teal!35!black,title={Before you move on}]
What does \hi{गुनगुनाना} mean? (\textquotedblleft{}To hum\textquotedblright{}.) Say it once more.
\end{tcolorbox}
